% SPDX-FileCopyrightText: © 2026 Idiap Research Institute <contact@idiap.ch>
%
% SPDX-FileContributor: Andre Freitas <andre.freitas@idiap.ch>
% Neuro-symbolic AI Group
%
% SPDX-FileContributor: Danilo Gusicuma <danilo.gusicuma@idiap.ch>
% Neuro-symbolic AI Group
%
% SPDX-License-Identifier: MIT
\documentclass[11pt]{article}
\usepackage[margin=1in]{geometry}
\usepackage[T1]{fontenc}
\usepackage[utf8]{inputenc}
\usepackage{lmodern}
\usepackage{amsmath,amssymb}
\usepackage{enumitem}
\usepackage{hyperref}
\title{Proof Discovery Language}
\author{Neuro-symbolic AI Group\\Idiap Research Institute}
\date{March 17, 2026}
\begin{document}
\maketitle
\begin{abstract}
PDL is a language for discovery programmes rather than only arguments of correctness. It treats coherence, conceptual nexus, material inference, obstruction, and governed recursion as first-class logical objects, and it distinguishes structural notation from the operations that materially transform one epistemic state into another.
\end{abstract}
\section{Core Claim}
The language is not centered on correctness-style proofs alone. It is centered on discovery programs: coherent nexuses of anchors, tensions, bridge objects, obstructions, and governed decisions over proof debt.
\section{Design Principles}
\begin{itemize}[leftmargin=*,itemsep=0.25em]
\item Prioritize conceptual nexus before closure.
\item Represent tensions and obstructions explicitly.
\item Treat bridge objects as typed inferential interfaces.
\item Make negative and impossibility reasoning first-class.
\item Make the underlying material-inference operators explicit rather than hiding them in prose.
\item Embed persistence, aggregation, and loops into the logic itself.
\end{itemize}
\section{Notation Layers}
\begin{itemize}[leftmargin=*,itemsep=0.25em]
\item \textbf{Structural Layer}: Names the epistemic objects under search: anchors, targets, tensions, nexuses, bridge objects, obstructions, programmes, debt, and gates. \emph{Binding.} This layer says what kind of thing is currently in play.
\item \textbf{Material-Inference Layer}: Names the operations that transform one epistemic object into another by licensed, non-monotonic, theorem-grounded moves. \emph{Binding.} This layer says what inferential move is being performed underneath the structural notation.
\item \textbf{Governance Layer}: Names persistence, aggregation, revision, continuation, and stop-control over discovery programmes. \emph{Binding.} This layer says how the process keeps state, loops, and changes direction without losing epistemic discipline.
\end{itemize}
\section{Primitives}
\begin{itemize}[leftmargin=*,itemsep=0.25em]
\item \textbf{Anchor} ($ A $): Grounded theorem body, definition, axiom, or imported result that licenses later moves. \emph{Material binding.} Material inference starts from anchors, not from topic labels. \emph{Platform binding.} KB statements and InformalMathLib entries.
\item \textbf{Target} ($ T $): The smallest nontrivial payoff or consequence worth organizing the search around. \emph{Material binding.} Discovery is guided by consequence pressure rather than free exploration. \emph{Platform binding.} Route target claims and queued conjectures.
\item \textbf{Tension} ($ X $): A typed gap showing why anchors do not yet reach the target: model mismatch, invariant weakness, transport loss, or barrier pressure. \emph{Material binding.} Tension is the primitive of discovery, because it identifies where reasoning must change form. \emph{Platform binding.} Statement deltas, obstruction classes, and failure diagnoses.
\item \textbf{Conceptual Nexus} ($ N $): A coherent local cluster of anchors, tensions, bridge objects, and barriers that can be reasoned over together. \emph{Material binding.} Discovery first asks whether there is a coherent nexus before demanding correctness-style closure. \emph{Platform binding.} Bridge portfolios, conceptual dynamics, and discovery branches.
\item \textbf{Bridge Object} ($ B $): A witness, invariant, certificate, transport law, normal form, or surrogate problem that could make the target reachable. \emph{Material binding.} Bridge objects express the conceptual breakthrough required by the search. \emph{Platform binding.} Bridge plans, abstraction conjectures, and candidate carriers.
\item \textbf{Obstruction} ($ O $): A barrier, counterexample pattern, impossibility theorem, or non-realization diagnosis that blocks a route. \emph{Material binding.} Negative mathematics is first-class discovery content, not only a failed proof attempt. \emph{Platform binding.} Blocked results, negative programs, and obstruction maps.
\item \textbf{Discovery Program} ($ P $): A paired positive/negative route over a nexus, together with proof debt, decision criteria, and revision triggers. \emph{Material binding.} A discovery program is the unit of guided scientific work, not an isolated proposition. \emph{Platform binding.} Proof sketches, queue entries, and branch dossiers.
\item \textbf{Proof Debt} ($ D $): The set of missing lemmas, unrealized carriers, transport laws, and closure obligations still required for completion. \emph{Material binding.} Debt is what makes discovery governable instead of open-ended. \emph{Platform binding.} Macro-strategy reports, proof obligations, and critique findings.
\item \textbf{Decision Gate} ($ G $): A criterion for continue, revise, import prior art, split regimes, or stop. \emph{Material binding.} Scientific discovery requires explicit transition rules, not only local generation. \emph{Platform binding.} Work loop, web-search gate, high-reward watch, and meta triggers.
\end{itemize}
\section{Material-Inference Operations}
\begin{itemize}[leftmargin=*,itemsep=0.25em]
\item \texttt{ext(A => c)}: Read a theorem body, definition, or axiom as a usable local commitment rather than as raw text. \emph{Effect.} Turns a grounded anchor into claim-bearing content. \emph{Universal binding.} licensed material consequence extraction \emph{Platform binding.} Phase 0 normalization, statement excerpts, theorem-body normalizer
\item \texttt{foc(A, T => X)}: Localize the exact tension between grounded anchors and the intended target. \emph{Effect.} Turns diffuse mismatch into a typed tension object. \emph{Universal binding.} problem localization / abductive targeting \emph{Platform binding.} statement deltas, failure diagnoses, obstruction classes
\item \texttt{conc(N => B)}: Turn a coherent conceptual nexus into a candidate bridge object, interface, or surrogate problem. \emph{Effect.} Moves from coherence pressure to an inspectable candidate object. \emph{Universal binding.} abductive bridge-object formation \emph{Platform binding.} abstraction conjectures, expert conjectures, branch-to-definition outputs
\item \texttt{tr\_I(B, s \textasciitilde{}> t)}: Move an invariant, witness condition, or theorem payload across a regime shift, representation change, or equivalence. \emph{Effect.} Makes cross-model or cross-language claims explicit and testable. \emph{Universal binding.} material transfer / invariant-preserving passage \emph{Platform binding.} transport lemmas, alignment work, HoTT/coherence routes
\item \texttt{fac(r => r1 o r2)}: Split one overloaded route into smaller theorem-sized interfaces or intermediate programmes. \emph{Effect.} Exposes hidden composition and missing middle objects. \emph{Universal binding.} decomposition / factorization reasoning \emph{Platform binding.} route decomposition, macro-strategy repair, categorical reframing
\item \texttt{comp(B1, B2 => B)}: Assemble several partial bridges, lemmas, or regime results into one route. \emph{Effect.} Promotes local inferential fragments into an end-to-end programme. \emph{Universal binding.} compositional synthesis \emph{Platform binding.} bridge portfolios, hybridization-of-partial-bridges, route composition
\item \texttt{obs(B => O)}: Derive the barrier, counterexample pattern, incoherence, or non-realization result that blocks a bridge. \emph{Effect.} Turns failed progress into explicit negative mathematics. \emph{Universal binding.} barrier derivation / impossibility localization \emph{Platform binding.} negative routes, blocked local theorems, contradiction-first probes
\item \texttt{rev(P, O => P')}: Update a discovery programme after an obstruction, failed proof step, or new prior-art import. \emph{Effect.} Makes conceptual change itself explicit rather than leaving it as narrative drift. \emph{Universal binding.} non-monotonic theory revision \emph{Platform binding.} revision loops, branch regeneration, meta-guided restarts
\item \texttt{agg(\{Ri\} => N*)}: Combine many local results, failed routes, or branch outcomes into a single higher-order nexus or search diagnosis. \emph{Effect.} Produces a global picture rather than isolated route state. \emph{Universal binding.} scientific synthesis / evidence aggregation \emph{Platform binding.} branch synthesis, breakthrough profiles, plausibility aggregation
\item \texttt{mem(R => M)}: Turn a proved, disproved, blocked, or imported result into durable problem memory. \emph{Effect.} Stops the system from rediscovering the same local mathematics repeatedly. \emph{Universal binding.} epistemic memory consolidation \emph{Platform binding.} InformalMathLib persistence, queue state, produced papers, prior-art imports
\item \texttt{route(P => next)}: Send a discovery programme to prove, further discovery, prior-art search, or stop. \emph{Effect.} Connects local reasoning to governance decisions. \emph{Universal binding.} controlled transition / research governance \emph{Platform binding.} work loop, meta triggers, high-reward continuation, queue promotion
\end{itemize}
\section{Judgments}
\begin{itemize}[leftmargin=*,itemsep=0.25em]
\item \texttt{K ; M |- A : grounded}: The KB plus InformalMathLib licenses anchor A as a grounded theorem-body object.
\item \texttt{A, X => N : coherent}: Anchors and tensions form a coherent conceptual nexus worth exploring.
\item \texttt{N \textasciitilde{}> B}: The nexus suggests a bridge object B as the best conceptual move.
\item \texttt{B [] I}: Bridge object B preserves invariant I across the intended route.
\item \texttt{B \textasciitilde{}> T}: If realized, bridge object B materially supports the target consequence T.
\item \texttt{B -| O}: Bridge object B is blocked or sharpened by obstruction O.
\item \texttt{N ; B |- P}: A coherent nexus and bridge object together induce a discovery program P.
\item \texttt{P |- D}: A discovery program exposes its proof debt explicitly.
\item \texttt{P ; O |- impossibility(T)}: A negative route yields a theorem-shaped impossibility or non-realization claim.
\item \texttt{P ; G @ next}: A decision gate routes the programme into prove, revision, prior-art import, or stop.
\end{itemize}
\section{Operational Judgments}
\begin{itemize}[leftmargin=*,itemsep=0.25em]
\item \texttt{A ; T |-m X}: Materially focus the mismatch between anchors and target into an explicit tension X.
\item \texttt{A ; X |-m N}: Stabilize anchors and tensions into a coherent conceptual nexus N.
\item \texttt{N |-m B}: Concretize the nexus into a bridge object, map, witness family, or surrogate problem B.
\item \texttt{B ; I ; s->t |-m tr\_I(B)}: Transport invariant I through bridge B from source regime s to target regime t.
\item \texttt{r |-m fac(r1,...,rn)}: Factor a route r into smaller theorem-sized route fragments.
\item \texttt{B |-m O}: Derive an obstruction O from the current bridge or route proposal.
\item \texttt{P ; O |-m P'}: Revise programme P into a new programme P' after a negative result or barrier localization.
\item \texttt{\{Ri\} |-m N*}: Aggregate many local route outcomes into a higher-order nexus or diagnosis N*.
\item \texttt{R |-m M}: Persist result R into durable memory object M.
\end{itemize}
\section{Coordination Logic}
\begin{itemize}[leftmargin=*,itemsep=0.25em]
\item \textbf{Coherence Before Closure}: A route is first judged as a coherent nexus of anchors, tensions, and bridge commitments before it is judged as correct.
\item \textbf{Positive/Negative Symmetry}: Every discovery program carries both a constructive route and a negative/impossibility route, each capable of yielding mathematical value.
\item \textbf{Material Inference First}: Inference moves are justified by theorem bodies, model regimes, barriers, and transport obligations before any correctness-style closure is attempted.
\item \textbf{Persistent Scientific Memory}: Queues, InformalMathLib, logs, and branch states are part of the logic itself because discovery depends on not forgetting failed and partial routes.
\item \textbf{Governed Recursion}: Loops are legitimate only when proof debt, obstruction pressure, or high-value subresults justify another cycle.
\end{itemize}
\section{Well-Formed Discovery Program}
\begin{itemize}[leftmargin=*,itemsep=0.25em]
\item At least one grounded anchor and one explicit target.
\item A typed tension explaining why the route is not immediate.
\item A coherent nexus rather than a bag of related statements.
\item A positive route, a negative route, and explicit proof debt.
\item A decision gate specifying continue, revise, import, split, or stop.
\end{itemize}
\section{Example Program}
\begin{itemize}[leftmargin=*,itemsep=0.25em]
\item A1: source theorem body
\item A2: target theorem body
\item A3: local barrier theorem
\item X1: transport from source regime to target regime loses the intended invariant
\item N1 := cohere(A1, A2, X1, barrier=A3)
\item B1 := witness family preserving the source-side hardness profile
\item ext(A1 => c1): extract the usable lower-bound content from the source theorem body
\item foc(A1, T1 => X1): isolate the exact transport loss between source and target
\item conc(N1 => B1): propose the witness family as the bridge object demanded by the nexus
\item tr\_I(B1, source \textasciitilde{}> target): test whether the invariant survives the intended transport
\item obs(B1 => O1): derive the non-realization barrier if the witness family collapses
\item rev(P1, O1 => P2): revise the route into a transport-obstruction programme
\item mem(O1 => M1): persist the blocked result into InformalMathLib as negative local theorem memory
\item If B1 is realizable and preserves the invariant, reduce target T to one route lemma L
\item If every realization of B1 collapses into the barrier, derive impossibility(B1 -> T)
\item \textbf{Debt.} L1: carrier realization; L2: transport theorem; L3: payoff forcing theorem
\item \textbf{Gate.} If L1 sharpens but L2 fails under barrier A3, reframe as transport-obstruction rather than continue constructive forcing
\end{itemize}
\end{document}
